\chapter*{Abstract}
\addcontentsline{toc}{chapter}{Abstract}

Earth Observation (EO) systems combine heterogeneous satellite metadata, raster assets, and
spatial queries whose reliability depends on both provider behavior and local processing
constraints. This thesis presents the design and implementation of a hybrid EO data
infrastructure for research and education in the Kvarken Space Center context. The system
combines public STAC and authenticated Copernicus Data Space Ecosystem (CDSE) access,
typed boundary validation, bounded asynchronous ingestion, content-addressed provenance,
SQLite-based spatial cataloguing, metadata quality profiling, and lightweight raster
analysis.

The implementation deliberately uses only the Python standard library at runtime. Injectable
transports and checked-in fixtures make failure recovery, OAuth2 token refresh, rate limiting,
concurrency, SQLite contention, and analytical behavior testable without network access. The
frozen empirical baseline ingests 256 synthetic scenes at concurrency 8 and records throughput,
payload footprint, catalog size, spatial query latency, quality distributions, and failure
recovery. The repository release baseline passes 53 offline tests, Ruff quality checks, and an
end-to-end health verification command.

\textbf{Keywords:} Earth Observation, STAC, CDSE, data infrastructure, provenance, SQLite,
bounded concurrency, reproducible research
